\documentclass[UTF8,a4paper,10pt]{ctexart}

\usepackage{amsmath}
\usepackage{geometry}
\usepackage{graphicx}
\usepackage{booktabs}
\usepackage{array}
\usepackage{float}
\usepackage{caption}
\usepackage{subcaption}
\usepackage{tabularx}
\usepackage{fancyhdr}
\usepackage{hyperref}
\usepackage{xcolor}

\geometry{left=1.55cm,right=1.55cm,top=1.45cm,bottom=1.55cm}
\hypersetup{colorlinks=true,linkcolor=black,urlcolor=black,citecolor=black}
\setlength{\parindent}{2em}
\setlength{\parskip}{2pt}
\setlength{\abovedisplayskip}{4pt}
\setlength{\belowdisplayskip}{4pt}
\setlength{\floatsep}{6pt}
\setlength{\textfloatsep}{6pt}
\setlength{\intextsep}{6pt}
\renewcommand{\arraystretch}{1.08}
\captionsetup{font=small,labelfont=bf}
\pagestyle{fancy}
\fancyhf{}
\fancyhead[L]{圆管在弯扭组合作用下的应力与内力测试实验}
\fancyhead[R]{\thepage}
\renewcommand{\headrulewidth}{0.3pt}

\newcommand{\ue}{\mu\varepsilon}
\newcommand{\img}[2]{\includegraphics[width=#1]{assets/#2}}

\begin{document}

\begin{center}
    {\Large\bfseries 圆管在弯扭组合作用下的应力与内力测试实验报告(由 Claude Opus 生成)}\\[4pt]
    {\normalsize 固体力学基础实验}\\[8pt]
    \begin{tabular}{>{\bfseries}r l >{\bfseries}r l >{\bfseries}r l >{\bfseries}r l}
        姓名： & 同学 & 学号： & 2024012843 & 班级： & 行健烽火4 & 日期： & 2026年5月26日
    \end{tabular}
\end{center}

\section{实验目的}

本实验以弯扭组合受力圆管为对象，利用电阻应变片和静态电阻应变仪测量圆管危险截面附近的表面应变。实验目的为：
\begin{enumerate}
    \item 掌握用电测法测定一点平面应力状态的基本方法；
    \item 根据弯曲与扭转的理论分析，选择合适的贴片位置和组桥方式；
    \item 通过应变测量分离弯矩与扭矩对应的应变分量；
    \item 根据重复加载数据判断测试结果的可靠性，并分析误差来源。
\end{enumerate}

\section{实验装置与原始资料}

实验装置由弯扭组合受力圆管、加载拐臂、电子万能试验机、DH3816N 静态应力应变测试分析系统、桥路接线箱和计算机采集软件组成。文件夹中共有 42 张 JPEG 图片，其中 21 张为原图，另外 21 张为重复副本；本报告只使用原图。另有《弯扭实验指导书.pdf》，用于确定实验名称、目的和报告整理要求。

\begin{figure}[H]
    \centering
    \begin{subfigure}[t]{0.48\textwidth}
        \centering
        \img{0.96\linewidth}{IMG_1866.jpg}
        \caption{DH3816N 静态应变测试系统与桥路接线}
    \end{subfigure}
    \hfill
    \begin{subfigure}[t]{0.48\textwidth}
        \centering
        \img{0.96\linewidth}{IMG_1867.jpg}
        \caption{圆管试样、贴片区与导线连接}
    \end{subfigure}
    \caption{实验设备与测点连接}
\end{figure}

\section{实验原理}

\subsection{电阻应变测量}

电阻应变片的电阻变化率与测点方向应变近似满足
\begin{equation}
    \frac{\Delta R}{R}=K\varepsilon ,
\end{equation}
其中 $K$ 为灵敏系数，$\varepsilon$ 为应变。应变仪通过惠斯通电桥测得桥路输出，并换算为微应变。单臂测量可以获得每个应变片的原始应变；半桥或全桥可以利用对称贴片的加减关系，放大目标内力分量并抵消温度、轴向拉压或弯曲等共同影响。

\subsection{圆管弯扭组合应力状态}

加载拐臂使空心圆管同时承受弯矩 $M$ 和扭矩 $T$。在圆管外表面危险点处，可近似视为平面应力状态：
\begin{equation}
    \sigma_x=\frac{MR}{I}, \qquad
    \tau_{x\theta}=\frac{TR}{J},
\end{equation}
其中 $R$ 为圆管外半径，$I$ 为截面对弯曲中性轴的惯性矩，$J$ 为极惯性矩。对空心圆管，
\begin{equation}
    I=\frac{\pi}{64}\left(D^4-d^4\right), \qquad
    J=\frac{\pi}{32}\left(D^4-d^4\right).
\end{equation}

材料处于线弹性阶段时，
\begin{equation}
    \sigma_x=E\varepsilon_x,\qquad
    \tau_{x\theta}=G\gamma_{x\theta},\qquad
    G=\frac{E}{2(1+\nu)} .
\end{equation}
主应力和主方向可由
\begin{equation}
    \sigma_{1,2}=\frac{\sigma_x}{2}\pm
    \sqrt{\left(\frac{\sigma_x}{2}\right)^2+\tau_{x\theta}^{\,2}},
    \qquad
    \tan 2\theta=\frac{2\tau_{x\theta}}{\sigma_x}
\end{equation}
求得。以上公式均未使用方框标注，便于直接排版和检查。

\subsection{本次数据的通道分离方法}

从三组 7 通道数据看，0-02 与 0-06 近似等大反号，适合作为弯曲对片；0-03 与 0-07 近似等大反号，适合作为 $\pm45^\circ$ 扭转对片。于是本报告按下式整理主要结果：
\begin{equation}
    \varepsilon_b=\frac{\varepsilon_{0-06}-\varepsilon_{0-02}}{2},
    \qquad
    \gamma_{x\theta}=\varepsilon_{0-07}-\varepsilon_{0-03}.
\end{equation}
其中 $\varepsilon_b$ 为圆管外表面弯曲应变幅值，$\gamma_{x\theta}$ 为由 $\pm45^\circ$ 片分离得到的工程剪应变。

\section{实验步骤}

\begin{enumerate}
    \item 检查圆管贴片区和接线端子，将各应变片接入 DH3816N 静态应变测试系统；
    \item 在 DHDAS 采集软件中建立 0-01 至 0-07 通道，完成清零和平衡；
    \item 通过电子万能试验机加载，照片中控制器读数显示试验力约为 $200\,\mathrm{N}$，负号代表加载方向；
    \item 保持载荷稳定后记录各通道微应变值，重复加载并保存三组有效数据；
    \item 对单臂测量数据进行整理，选择对称通道计算弯曲应变、剪应变、主应力和主方向。
\end{enumerate}

\section{原始数据整理}

\subsection{加载读数}

控制器照片显示加载过程中试验力基本稳定在 $200\,\mathrm{N}$ 左右，位置读数约为 $-291$ 至 $-295\,\mathrm{mm}$。表中负号只表示试验机约定的方向，后续分析采用载荷大小。

\begin{table}[H]
    \centering
    \caption{电子万能试验机控制器读数}
    \begin{tabular}{lccc}
        \toprule
        照片 & 试验力/N & 位置/mm & 位移/mm \\
        \midrule
        IMG\_1869 & $-198.5$ & $-291.246$ & $0.0008$ \\
        IMG\_1870 & $-201.2$ & $-291.246$ & $0.0008$ \\
        IMG\_1871 & $-195.4$ & $-295.056$ & $0.0010$ \\
        \bottomrule
    \end{tabular}
\end{table}

\begin{figure}[H]
    \centering
    \begin{subfigure}[t]{0.32\textwidth}
        \centering
        \img{0.98\linewidth}{IMG_1869.jpg}
        \caption{约 $201\,\mathrm{N}$}
    \end{subfigure}
    \hfill
    \begin{subfigure}[t]{0.32\textwidth}
        \centering
        \img{0.98\linewidth}{IMG_1870.jpg}
        \caption{约 $204\,\mathrm{N}$}
    \end{subfigure}
    \hfill
    \begin{subfigure}[t]{0.32\textwidth}
        \centering
        \img{0.98\linewidth}{IMG_1871.jpg}
        \caption{约 $197\,\mathrm{N}$}
    \end{subfigure}
    \caption{加载控制器照片}
\end{figure}

\subsection{七通道应变读数}

三次有效采集的读数如表 \ref{tab:strain-data} 所示，单位均为 $\ue$。读数来自采集软件表格中高亮或稳定行。

\begin{table}[H]
    \centering
    \caption{7 通道单臂应变测量结果}
    \label{tab:strain-data}
    \resizebox{\textwidth}{!}{
    \begin{tabular}{lrrrrrrr}
        \toprule
        照片 & 0-01 & 0-02 & 0-03 & 0-04 & 0-05 & 0-06 & 0-07 \\
        \midrule
        IMG\_1868 & 27.4 & -273.1 & -215.7 & 73.7 & -24.2 & 278.8 & 214.8 \\
        IMG\_1872 & 25.9 & -270.9 & -215.9 & 70.8 & -24.9 & 276.6 & 212.4 \\
        IMG\_1876 & 37.9 & -254.8 & -201.6 & 84.2 & -12.2 & 286.3 & 222.3 \\
        \midrule
        平均值 & 30.4 & -266.3 & -211.1 & 76.7 & -20.4 & 280.6 & 216.5 \\
        标准差 & 6.5 & 10.0 & 8.2 & 6.5 & 7.1 & 5.1 & 5.2 \\
        \bottomrule
    \end{tabular}}
\end{table}

\begin{figure}[H]
    \centering
    \begin{subfigure}[t]{0.32\textwidth}
        \centering
        \img{0.98\linewidth}{IMG_1868.jpg}
        \caption{第一次采集}
    \end{subfigure}
    \hfill
    \begin{subfigure}[t]{0.32\textwidth}
        \centering
        \img{0.98\linewidth}{IMG_1872.jpg}
        \caption{第二次采集}
    \end{subfigure}
    \hfill
    \begin{subfigure}[t]{0.32\textwidth}
        \centering
        \img{0.98\linewidth}{IMG_1876.jpg}
        \caption{第三次采集}
    \end{subfigure}
    \caption{DHDAS 软件中的 7 通道应变显示}
\end{figure}

\subsection{单通道补充记录}

部分照片只显示 0-01 单通道，用于说明清零、单片读取和调试过程。由于其工况与 7 通道重复加载记录不完全对应，未并入主计算。

\begin{table}[H]
    \centering
    \caption{0-01 单通道补充照片读数}
    \begin{tabular}{lcc}
        \toprule
        照片 & 0-01 读数/$\ue$ & 说明 \\
        \midrule
        IMG\_1878 & 0.8 & 清零或近零状态 \\
        IMG\_1880 & 242.1 & 单通道加载读数 \\
        IMG\_1883 & 22.1 & 小应变调试读数 \\
        IMG\_1888 & -188.9 & 单通道反向读数 \\
        \bottomrule
    \end{tabular}
\end{table}

\section{数据处理与计算}

\subsection{弯曲应变与剪应变分离}

由 0-02/0-06 对片和 0-03/0-07 对片计算得到表 \ref{tab:derived-strain}。

\begin{table}[H]
    \centering
    \caption{分离得到的弯曲应变和剪应变}
    \label{tab:derived-strain}
    \begin{tabular}{lcc}
        \toprule
        照片 & $\varepsilon_b=(\varepsilon_{0-06}-\varepsilon_{0-02})/2$ / $\ue$
             & $\gamma_{x\theta}=\varepsilon_{0-07}-\varepsilon_{0-03}$ / $\ue$ \\
        \midrule
        IMG\_1868 & 276.0 & 430.5 \\
        IMG\_1872 & 273.8 & 428.3 \\
        IMG\_1876 & 270.6 & 423.9 \\
        \midrule
        平均值 & 268.2 & 415.2 \\
        标准差 & 2.7 & 3.4 \\
        变异系数 & 0.99\% & 0.79\% \\
        \bottomrule
    \end{tabular}
\end{table}

可见，对称组桥后得到的弯曲应变和剪应变重复性明显优于单个通道，三次重复加载的离散程度小于 1\%，满足实验指导书中对重复测量稳定性的要求。

\subsection{应力与主应力计算}

按钢材常用弹性参数 $E=206\,\mathrm{GPa}$、$\nu=0.30$ 计算，剪切模量为
\begin{equation}
    G=\frac{206}{2(1+0.30)}=79.2\,\mathrm{GPa}.
\end{equation}

将平均弯曲应变和剪应变代入本构关系：
\begin{align}
    \sigma_x &= E\varepsilon_b
    =206000\times 268.2\times10^{-6}
    =55.2\,\mathrm{MPa},\\
    \tau_{x\theta} &= G\gamma_{x\theta}
    =79230.8\times 415.2\times10^{-6}
    =32.9\,\mathrm{MPa}.
\end{align}

危险点处主应力为
\begin{align}
    \sigma_1
    &=\frac{55.2}{2}+\sqrt{\left(\frac{55.2}{2}\right)^2+32.9^2}
    =70.8\,\mathrm{MPa},\\
    \sigma_2
    &=\frac{55.2}{2}-\sqrt{\left(\frac{55.2}{2}\right)^2+32.9^2}
    =-15.6\,\mathrm{MPa}.
\end{align}
主方向相对圆管轴向的夹角为
\begin{equation}
    \theta=\frac12\arctan\left(\frac{2\times32.9}{55.2}\right)
    =25.0^\circ .
\end{equation}

\subsection{内力反算关系}

若补充圆管外径 $D$、内径 $d$、弯矩力臂 $a$ 和扭矩力臂 $b$，可进一步计算实验反算内力并与理论值比较：
\begin{equation}
    M_{\mathrm{exp}}=\frac{\sigma_x I}{R},\qquad
    T_{\mathrm{exp}}=\frac{\tau_{x\theta} J}{R},
\end{equation}
\begin{equation}
    M_{\mathrm{th}}=Pa,\qquad T_{\mathrm{th}}=Pb.
\end{equation}
本次照片没有记录 $D,d,a,b$ 的实测数值，因此本报告不编造理论内力和相对误差。现有照片数据足以完成应变分离、重复性检验、应力状态计算和主方向计算；若后续补充几何尺寸，可直接按上式补全内力对比表。

\section{结果分析}

第一，0-02 与 0-06 通道、0-03 与 0-07 通道均呈现近似等大反号关系，说明贴片和桥路分组符合弯扭组合圆管的对称性。0-02/0-06 的平均幅值差为
\begin{equation}
    \frac{280.6-266.3}{(280.6+266.3)/2}=5.2\%,
\end{equation}
0-03/0-07 的平均幅值差为
\begin{equation}
    \frac{216.5-211.1}{(216.5+211.1)/2}=2.5\%.
\end{equation}
两组主用通道的对称性较好，其中扭转对片更稳定。

第二，直接单通道数据中 0-01、0-04、0-05 的离散性较大，主要原因可能是这些通道位于应变较小区域、接线或零点漂移影响占比更高。采用对称通道做差后，共同漂移被显著抵消，因此分离出的 $\varepsilon_b$ 和 $\gamma_{x\theta}$ 重复性更高。

第三，主应力计算表明该危险点处为拉压主应力并存状态，最大主应力约 $70.8\,\mathrm{MPa}$，最小主应力约 $-15.6\,\mathrm{MPa}$，主方向与圆管轴线夹角约 $25.0^\circ$。这与弯曲正应力和扭转剪应力共同作用下的平面应力状态一致。

\section{误差来源}

实验误差主要来自以下方面：
\begin{enumerate}
    \item 贴片位置和方向误差：应变片若未严格沿轴向或 $\pm45^\circ$ 方向粘贴，会影响应变分离；
    \item 桥路接线误差：端子接触电阻、导线编号混淆或桥臂选择不当会改变读数符号和幅值；
    \item 零点漂移：应变仪清零后仍可能随时间、温度和导线状态发生小幅漂移；
    \item 加载稳定性：控制器显示载荷约在 $197$ 至 $204\,\mathrm{N}$ 之间波动，会造成应变读数轻微变化；
    \item 几何参数缺失：本次原始照片未记录圆管外径、内径和力臂长度，限制了内力理论值的数值比较。
\end{enumerate}

\section{结论}

本实验完成了弯扭组合圆管在约 $200\,\mathrm{N}$ 加载下的静态应变测量。根据三组 7 通道数据，0-02/0-06 可用于分离弯曲应变，0-03/0-07 可用于分离扭转剪应变。处理结果为
\[
    \varepsilon_b=268.2\,\ue,\qquad
    \gamma_{x\theta}=415.2\,\ue .
\]
按 $E=206\,\mathrm{GPa}$、$\nu=0.30$ 计算，危险点处
\[
    \sigma_x=55.2\,\mathrm{MPa},\qquad
    \tau_{x\theta}=32.9\,\mathrm{MPa},
\]
主应力为
\[
    \sigma_1=70.8\,\mathrm{MPa},\qquad
    \sigma_2=-15.6\,\mathrm{MPa},\qquad
    \theta=25.0^\circ .
\]
对称通道分离后的重复性优于 1\%，说明本次测量数据稳定，弯扭组合应力状态判断合理。

\appendix
\section{原始照片索引}

下图为 21 张原始照片的汇总索引，用于说明本报告所依据的现场记录范围。

\begin{figure}[H]
    \centering
    \includegraphics[height=0.86\textheight]{assets/contact_sheet.jpg}
    \caption{原始照片汇总索引}
\end{figure}

\end{document}
